\label{sec:formulation}

\subsection{CST parameterisation and its linearity}
\label{sec:cst-linearity}
% <- T2; math source \dossierref{\S3.1}--\S3.4

Let $\psi = x/c \in [0,1]$ and $\zeta = z/c$. The CST class function for a
round-nose, sharp-tail family is
\begin{equation}
  C_{N_1}^{N_2}(\psi) = \psi^{N_1}(1-\psi)^{N_2}, \qquad N_1 = 0.5,\; N_2 = 1.0,
\end{equation}
and the shape function is a Bernstein-polynomial expansion of order $n$ per
surface,
\begin{equation}
  S_i(\psi) = K_i\,\psi^i(1-\psi)^{n-i}, \quad K_i = \binom{n}{i}, \qquad
  S_{u}(\psi) = \sum_{i=0}^{n} A_{u,i} S_i(\psi), \quad
  S_{l}(\psi) = \sum_{i=0}^{n} A_{l,i} S_i(\psi).
\end{equation}
The upper and lower surfaces are
\begin{equation}
  \zeta_{u}(\psi) = C(\psi)\,S_u(\psi) + \psi\,\zeta_{T,u}, \qquad
  \zeta_{l}(\psi) = C(\psi)\,S_l(\psi) + \psi\,\zeta_{T,l},
  \label{eq:surfaces}
\end{equation}
with $\zeta_{T,u} = +t_{TE}/2$, $\zeta_{T,l} = -t_{TE}/2$. The sign
convention is binding throughout: lower-surface coefficients are stored with
their natural negative sign, so the shared-leading-edge row reads
$A_{u,0}+A_{l,0}=0$.

Figure~\ref{fig:cst-basis} shows the construction term by term. The class
function of panel~(a) is fixed geometry that carries no dependence on $A$ at
all; the Bernstein shape functions of panel~(b) are fixed polynomials; and the
surface of panel~(d) is the $A$-weighted sum of the terms in panel~(c). Only
the choice of $N_1$ and $N_2$ is nonlinear, and those are held fixed
throughout this work at the round-nosed, sharp-trailing-edge values
$N_1=\tfrac12$, $N_2=1$.

\begin{figure}[htbp]
\centering
\includegraphics[width=\textwidth]{fig_cst_basis.png}
\caption{The CST construction, evaluated rather than sketched.
(a)~the class function for three $(N_1,N_2)$ pairs, with the airfoil-like
choice used here highlighted; (b)~the Bernstein shape functions at $n=8$;
(c)~the weighted terms $A_i C(\psi) S_i(\psi)$ for a fitted NACA 2412 upper
surface, and their sum; (d)~the resulting surface against the section it was
fitted to. Every panel except the class function is linear in $A$, which is
the property Eq.~\eqref{eq:dsurface-dA} states and the rest of the
formulation depends on.}
\label{fig:cst-basis}
\end{figure}

The property this paper exploits is that Eq.~\eqref{eq:surfaces} is
\emph{linear} in $A$: the partial derivative
\begin{equation}
  \frac{\partial \zeta}{\partial A_i} = C(\psi)\,S_i(\psi)
  \label{eq:dsurface-dA}
\end{equation}
does not depend on $A$. It is a fixed matrix, computed once from the
chordwise node distribution and cached (T2 requirement,
\dossierref{\S7.3}): \emph{``does not
depend on A. Recomputing it per Newton iteration is the single most common
way to lose the performance argument.''} This design-independence is what
makes the $\partial R/\partial A$ Jacobian block of \S\ref{sec:newton-system}
assemble as a cheap matrix product rather than a re-derived, geometry-
dependent operator at every Newton step.

Two endpoint identities follow directly from Eq.~\eqref{eq:surfaces} and are
used as constraint rows in \S\ref{sec:constraints}. The trailing-edge
identity is $\zeta'(1) = \zeta_T - A_n$ at $N_2=1$, and it does \emph{not}
take the same form on both surfaces: with the boat-tail half-angle defined
per surface as the wedge closing angle, the upper surface descends toward the
trailing edge ($\tan\beta_u = -\zeta_u'(1)$) while the lower surface rises
($\tan\beta_l = +\zeta_l'(1)$), so
\begin{equation}
  A_0 = \sqrt{2R_{LE}/c}, \qquad
  A_{u,n} = \zeta_{T,u}/c + \tan\beta_u, \qquad
  A_{l,n} = \zeta_{T,l}/c - \tan\beta_l.
  \label{eq:endpoint}
\end{equation}
The sign flip on the lower surface is stated explicitly because applying the
upper-surface form to a lower-surface coefficient, which is negative under
the stored-sign convention, silently produces the wrong constraint row; that
error occurred once during development and is now pinned by a test against
the fitted trailing-edge slope. Near $\psi\to0$,
$\zeta \approx A_0\sqrt{\psi}$, so $R_{LE} = A_0^2/2$; $G^2$ curvature
continuity across the leading edge reduces to matching the leading
coefficient magnitude on the two surfaces, $A_{u,0} = |A_{l,0}|$; under the
stored-sign convention ($A_{l,i}<0$), the row $A_{u,0} + A_{l,0} = 0$.

Figure~\ref{fig:cst-parameters} gives these identities their physical
reading and checks the design-independence claim numerically. Panels~(a)
and~(b) plot the two endpoint identities of Eq.~\eqref{eq:endpoint} with the
fitted NACA 2412 values marked, so the coefficients that the inverse solve
returns can be read directly as a nose radius and a trailing-edge wedge angle
rather than as abstract weights. Panel~(c) is the one that matters for the
Newton system: the markers are the surface displacement measured by actually
perturbing one coefficient of the fitted section and re-evaluating the
geometry, and the lines are the analytic prediction $\delta A_i\,C S_i$. They
agree to the level printed on the panel, and the agreement does not depend on
where the remaining coefficients sit, which is exactly what permits
$\partial\zeta/\partial A$ to be assembled once and cached.

\begin{figure}[htbp]
\centering
\includegraphics[width=\textwidth]{fig_cst_parameters.png}
\caption{What the coefficients mean, and why the geometric Jacobian block is
constant. (a)~leading-edge radius against $A_0$, from
Eq.~\eqref{eq:endpoint}, with the fitted NACA 2412 value marked;
(b)~trailing-edge wedge half-angle against the final coefficient, for both
surfaces under the stored-sign convention; (c)~superposition checked rather
than asserted, markers measured by perturbing a single coefficient and lines
predicted analytically; (d)~the fitted coefficient vector itself, upper and
lower blocks, in the ordering used throughout
(\S\ref{sec:newton-system}).}
\label{fig:cst-parameters}
\end{figure}

\subsection{Constraints as linear rows}
\label{sec:constraints}
% <- T3; closed-form source \dossierref{\S3.4}

Because Eq.~\eqref{eq:surfaces} is linear in $A$, every geometric functional
that is itself linear or integral in $\zeta$ becomes a linear row
$g\cdot A = b$ rather than a nonlinear predicate evaluated on a rebuilt
mesh. The leading-edge-radius and trailing-edge-wedge rows follow
immediately from Eq.~\eqref{eq:endpoint}. The one requiring a closed-form
integral is inscribed area: the per-surface contribution is
\begin{equation}
  \int_0^1 C(\psi)S_i(\psi)\,d\psi
    = K_i \int_0^1 \psi^{i+N_1}(1-\psi)^{n-i+N_2}\,d\psi
    = K_i\,B(i+N_1+1,\, n-i+N_2+1),
  \label{eq:area-beta}
\end{equation}
a Beta function, plus the trailing-edge term $\int_0^1 \psi\,\zeta_T\,d\psi =
\zeta_T/2$. Eq.~\eqref{eq:area-beta} is evaluated with
\texttt{scipy.special.beta}, never by numerical quadrature in a production
constraint row; quadrature is retained only as an independent gate check. That gate passed: \texttt{area\_row} matches
numerical quadrature of the fitted NACA~2412 to $10^{-10}$. This closed-form machinery
is the same Beta-function apparatus Joseph \& Mohan~\citep{joseph2021closed}
use for thin-airfoil lift/moment integrals against a Bernstein camber basis.
This is not a coincidence: both integrate Bernstein polynomials against the
class function, and it is the analytic kernel reused for the linear pre-solve
of \S\ref{sec:presolve}.

Under a B-spline-with-wrapper representation, by contrast, every row in
Table~\ref{tab:constraint-comparison} is a nonlinear or non-smooth predicate
evaluated on a rebuilt mesh. This is the stack that forced surrogate-based
search in the author's earlier turbine-blading
work~\citep{sathish2019sco2} (\dossierref{\S1.5}).

\paragraph{Which rows the reported results actually exercise.} The
constraint machinery described here is broader than what the experiments in
\S\ref{sec:results} use, and the difference should be stated rather than left
to be inferred. Every configuration reported in this paper runs either with
the leading edge prescribed, in which case $K=0$ and no constraint row is
assembled at all, or with \texttt{le\_treatment=none}, in which case the
single shared-leading-edge row $A_{u,0}+A_{l,0}=0$ is used. The
trailing-edge-wedge and inscribed-area rows, including the Beta-function
derivation of Appendix~\ref{app:area}, are implemented and unit-tested
against independent quadrature but are not exercised by any result reported
here. They are presented as part of the formulation because they follow from
the same linearity and cost nothing to assemble, not as validated
contributions to the measured outcomes.

\begin{table}[h]
\centering
\caption{Constraint representation under CST against a B-spline
control-point wrapper. Every row that is a linear, exact algebraic condition
on the CST coefficients becomes a nonlinear or non-smooth predicate on a
rebuilt mesh under the wrapper, which is what forces an outer search rather
than a root-find. Source: \dossierref{\S1.5}.}
\label{tab:constraint-comparison}
\footnotesize
\begin{tabularx}{\textwidth}{@{}l L L@{}}
\toprule
Constraint & Under CST & Under the B-spline wrapper \\
\midrule
LE radius & $A_0=\sqrt{2R_{LE}/c}$: linear, exact & implicit, via tangent points \\
TE thickness & $\zeta_T$ term: linear & geometric post-check \\
Boat-tail/wedge angle & $A_n$: linear & nonlinear angular DV \\
Inscribed area & linear functional, Beta-function closed form & numerical, non-smooth \\
Curvature & rational function of $A$, differentiable & monotonicity predicate, non-smooth \\
\bottomrule
\end{tabularx}
\end{table}

\subsection{The extended Newton system and DOF accounting}
\label{sec:newton-system}
% <- T5, FM-1-\S5, dossier \S7.6

The extended system appends the CST coefficients (and, optionally, angle of
attack) directly to the flow solver's own global unknown vector. Unknowns:
\begin{equation}
  U \;(\text{flow state}, N_{tot}), \qquad
  A \;(\text{CST coefficients}, n_A), \qquad
  \alpha \;(\text{angle of attack}, 0\text{ or }1).
\end{equation}
Equations:
\begin{equation}
  R(U; x(A)) = 0 \;\;(N_{tot}), \qquad
  T(U) - Cp_{target} = 0 \;\;(M), \qquad
  G\cdot A - b = 0 \;\;(K),
\end{equation}
with the block Jacobian, dimensions annotated below each block for the T7
winning configuration. There, mfoil's global system carries
$N_{sys}=230$ stations (200 airfoil-surface nodes plus wake,
\S\ref{sec:falsifiable}) with four boundary-layer state variables each, so
$N_{tot}=4N_{sys}=920$; the prescribed-LE treatment leaves
$n_{A,\text{free}}=16$ of the 18 CST coefficients as unknowns; and $M=16$,
$K=0$, $n_\alpha=0$. The assembled system is therefore
$J\in\mathbb{R}^{936\times936}$, and the recorded rank at the first iteration
is 936, that is, full
(\pth{experiments/results/t7_naca2412/diagnostics.json}):
\begin{equation}
  J = \left[
  \begin{array}{c|c|c}
    \displaystyle\underbrace{\partial R/\partial U}_{N_{tot}\times N_{tot}} &
    \displaystyle\underbrace{\partial R/\partial A}_{N_{tot}\times n_{A,\text{free}}} &
    \displaystyle\underbrace{\partial R/\partial \alpha}_{N_{tot}\times n_\alpha} \\[1.1em] \hline
    \displaystyle\underbrace{\partial T/\partial U}_{M\times N_{tot}} &
    \displaystyle\underbrace{0}_{M\times n_{A,\text{free}}} &
    \displaystyle\underbrace{\partial T/\partial \alpha}_{M\times n_\alpha} \\[1.1em] \hline
    \displaystyle\underbrace{0}_{K\times N_{tot}} &
    \displaystyle\underbrace{\partial G/\partial A}_{K\times n_{A,\text{free}}} &
    \displaystyle\underbrace{0}_{K\times n_\alpha}
  \end{array}
  \right]
  \;\in\; \mathbb{R}^{(N_{tot}+M+K)\times(N_{tot}+n_{A,\text{free}}+n_\alpha)},
  \label{eq:jacobian-blocks}
\end{equation}
square by Eq.~\eqref{eq:dof}'s construction-time assertion
($M+K=n_{A,\text{free}}+n_\alpha$), so the $M+K$ extra rows exactly balance
the $n_{A,\text{free}}+n_\alpha$ extra columns beyond the
$N_{tot}\times N_{tot}$ flow block, which is itself always square. The
$M\times n_{A,\text{free}}$ zero block in the middle column reflects
that the target-Cp equations $T(U)-Cp_{target}$ depend on $A$ only
\emph{through} $U$ (Cp is read off the converged flow state, not directly
from geometry). Its entire $A$-sensitivity is therefore carried through the
$\partial R/\partial A$ column and the chain of Newton updates; there is
no direct $T$-vs-$A$ coupling to assemble. $\partial R/\partial U$ and
$\partial R/\partial\alpha$ are mfoil's own
analytic, sparsely-assembled blocks, reused unchanged. $\partial T/\partial
U$ is a near-selection matrix (Cp at a node as a direct function of the
panel/edge-velocity unknowns). $\partial G/\partial A = G$ is constant and
free, per \S\ref{sec:constraints}. $\partial R/\partial A =
(\partial R/\partial x)(\partial x/\partial A)$ is the block that carries the
architectural claim; its second factor is the cached, design-independent
matrix of Eq.~\eqref{eq:dsurface-dA}, and its first factor is addressed in
\S\ref{sec:derivatives}.

\paragraph{DOF accounting.} The dossier's original bookkeeping
(\dossierref{\S4}~FM-1, \S7.6) is
\[
  \mathrm{DOF} = 2(n+1) - 1_{\text{shared LE}} - 1_{\text{TE fixed}}
  + 3_{\text{scale/rotate/translate}} + 1_{\text{stagger or }\alpha}
  \quad\stackrel{!}{=}\quad M + K.
\]
In the Stage-1 (isolated-airfoil, mfoil) implementation, the three rigid
overall modes (scale/rotate/translate) are not exposed as free Newton
unknowns, so the implemented squareness condition is the simpler, equivalent
statement actually asserted at construction time and logged every run:
\begin{equation}
  M + K = n_{A,\text{free}} + n_{\alpha},
  \label{eq:dof}
\end{equation}
where $n_{A,\text{free}}$ is the count of CST coefficients \emph{not}
eliminated by a prescribed-LE treatment (\S\ref{sec:falsifiable}) or an
equality-row elimination, and $n_\alpha\in\{0,1\}$ depending on whether angle
of attack is a free unknown for the run in question. Violating
Eq.~\eqref{eq:dof} raises at construction, before any Newton step. This is
the FM-1 guard (\S\ref{sec:fm}), and it is exercised directly in the T8
ablation matrix (\S\ref{sec:ablations}): deliberately over- and under-
squaring the baseline system (n=8/side, $n_{A,\text{free}}=16$, $n_\alpha=0$)
produces the exact, actionable diagnostics
\begin{quote}\footnotesize\ttfamily
"extended system not square: M+K=17 but\\
\hspace*{1em}n\_A\_free+n\_alpha=16 (M=17, K=0, n\_A\_free=16)"\\[2pt]
"extended system not square: M+K=15 but\\
\hspace*{1em}n\_A\_free+n\_alpha=16 (M=15, K=0, n\_A\_free=16)"
\end{quote}
with zero Newton iterations spent and no silent wrong answer. Getting Eq.~\eqref{eq:dof} wrong is,
per the dossier, plausibly what the author's PhD-era divergence was
misdiagnosed as an LE-singularity problem when it was rank deficiency
(\dossierref{\S4}, FM-1). Instrumenting
for it directly, rather than inferring it from Newton behaviour, is the
point of Eq.~\eqref{eq:dof} being an assertion rather than a derived
quantity.

\subsection{Derivative strategy: why finite-difference-over-\texorpdfstring{$A$}{A}, not the analytic chain rule or complex-step}
\label{sec:derivatives}
% <- derivative-mode decision tree; empirical finding (T1)

The derivative-mode decision tree adopted for this work prefers, in order:
(1) the analytic chain rule $\partial R/\partial A = (\partial R/\partial x)(\partial
x/\partial A)$ if mfoil exposes $\partial R/\partial x$ analytically; (2)
complex-step over $A$ if the residual-assembly path is free of
\texttt{abs}/\texttt{max}/\texttt{min}/comparison operators; (3) real
(central) finite differences over $A$ otherwise. T1 introspection resolved both questions empirically, and
neither resolved in the analytic-chain-rule's favour:

\begin{itemize}
\item \textbf{The analytic chain rule is incomplete for this solver class.}
  mfoil exposes \texttt{M.glob.R\_x}, but it is only a \emph{partial}
  sensitivity: real and useful for the arc-length ($\xi$) mapping, but
  missing the geometry-dependent operators that matter most, the
  inviscid panel/AIC coupling and the stagnation-point formulation. Using this partial
  block as if it were the complete $\partial R/\partial x$ would silently
  drop exactly the rows the leading-edge failure mode (FM-3,
  \S\ref{sec:fm}) is about.
\item \textbf{Complex-step is blocked, not merely inconvenient.} A bare
  \texttt{residual\_station}\slash\allowbreak\texttt{residual\_transition} call, with fixed
  panels and a complex-perturbed state, is verified complex-step-safe and
  reproduces the analytic $R_x$ to float64 machine epsilon. The \emph{full} augmented residual,
  which is what a perturbation in $A$ actually varies since $A$ moves
  the panel geometry itself, calls the panel-influence and
  stagnation-detection code paths. These use \texttt{atan2}, branch
  selection on real quantities, and a hard \texttt{ValueError} guard for
  complex node arrays. These are exactly the operators complex-step
  cannot tolerate, and they sit in precisely the geometry-dependent block
  that (1) shows was missing from the analytic path: complex-step
  \emph{cannot fill the gap the analytic chain rule leaves}.
\end{itemize}

The implemented and gate-verified strategy is therefore real central finite
differences with respect to $A$ only, never with respect to node
coordinates $x$ ().
This is a deliberate architectural discipline, not a fallback of
convenience. $A$ has $n_A\approx16$--$36$ columns depending on Bernstein
order, so a full central-difference Jacobian block costs
$O(n_A)$ residual \emph{evaluations} per Newton step, not flow solves,
against the alternative of finite-differencing over $N_{panel}\gg n_A$ node
coordinates. That alternative would silently destroy the flow-solve-count
argument this paper's Results section depends on. T5 gate closure verified
the resulting Jacobian columns to 5--6 significant figures across a
step-size sweep with no stencil discontinuities. The step size itself ($\mathrm{fd\_step}=10^{-6}$,
\pth{configs/default.yaml}) is not hand-tuned per case. The T5 sweep
varied it across several decades and confirmed the resulting FD-column
agreement with the (independently verified, \S\ref{sec:derivatives} above)
complex-step-safe partial block was flat across the sweep's central range,
with no visible truncation/roundoff trade-off cliff at the scales exercised. That is, $10^{-6}$ sits well inside a
broad plateau, not at a knife-edge optimum a different problem instance
would move.

\subsection{Algorithmic summary}
\label{sec:algorithms}

Figure~\ref{fig:pipeline} gives the whole procedure in one view before the
pseudocode states it precisely. The two shaded regions are the paper's two
structural claims: everything left of the Newton loop is set up once and
never revisited, and the loop itself solves a single square system rather
than driving an outer optimiser. The two decision points are the only places
the procedure can refuse to proceed, and both are diagnostics rather than
tolerances: the squareness assertion of Eq.~\eqref{eq:dof} is a statement
about degrees of freedom, and the realisability gate is a statement about
whether the target lies in the reachable set at all.

\begin{figure}[htbp]
\centering
\begin{tikzpicture}[node distance=6.5mm and 9mm]
  \node[fcstart] (tgt) {Target $C_p$ or $u_e$\\ \scriptsize (generated, or user-supplied)};
  \node[fcbox, below=of tgt] (fit) {Fit baseline to CST\\ \scriptsize $A_0$, $\zeta_T$, order $n$};
  \node[fcbox, below=of fit] (pre) {T4 linear pre-solve\\ \scriptsize sensitivity matrix $M$, KKT initial guess};
  \node[fcdec, below=8mm of pre] (real) {realisable?};
  \node[fcfail, right=14mm of real] (warn) {report\\ realisability};
  \node[fcbox, below=8mm of real] (qr) {QR-pivoted station selection\\ \scriptsize restricted to $x \geq$ prescribed-LE fraction};
  \node[fcdec, below=8mm of qr] (sq) {$M{+}K = n_{A,\mathrm{free}}{+}n_\alpha$?};
  \node[fcfail, right=14mm of sq] (dof) {DOF error,\\ zero iterations};

  \node[fcbox, below=8mm of sq] (newt) {\textbf{Extended Newton step}\\
        \scriptsize flow residual $R$, target rows $T$, constraint rows $GA-b$\\
        \scriptsize sparse LU on one square system};
  \node[fcdec, below=8mm of newt] (conv) {converged?};
  \node[fcout, below=8mm of conv] (rel) {Release transition,\\ direct-solve, verify};
  \node[fcout, below=of rel] (out) {Recovered geometry $A$,\\ \scriptsize diagnostics D-1 to D-6};

  \begin{scope}[on background layer]
    \node[fit=(tgt)(qr), draw=NavyBlue!35, fill=NavyBlue!5, rounded corners,
          inner sep=3.2mm] (setup) {};
    \node[fit=(newt)(conv), draw=ForestGreen!40, fill=ForestGreen!6,
          rounded corners, inner sep=3.2mm] (loop) {};
  \end{scope}
  \node[anchor=south west, font=\scriptsize\itshape, NavyBlue!70]
        at ([yshift=0.8mm]setup.north west) {set up once, before any iteration};
  \node[anchor=north west, font=\scriptsize\itshape, ForestGreen!60]
        at ([yshift=-0.8mm]loop.south west) {square system, no outer loop};

  \draw[fcarrow] (tgt) -- (fit);
  \draw[fcarrow] (fit) -- (pre);
  \draw[fcarrow] (pre) -- (real);
  \draw[fcarrow] (real) -- node[fclab,above] {no} (warn);
  \draw[fcarrow] (real) -- node[fclab,right] {yes} (qr);
  \draw[fcarrow] (qr) -- (sq);
  \draw[fcarrow] (sq) -- node[fclab,above] {no} (dof);
  \draw[fcarrow] (sq) -- node[fclab,right] {yes} (newt);
  \draw[fcarrow] (newt) -- (conv);
  \draw[fcarrow] (conv) -- node[fclab,right] {yes} (rel);
  \draw[fcarrow] (rel) -- (out);
  \draw[fcarrow] (conv.east) -- ++(9mm,0) node[fclab,midway,above] {no}
        |- (newt.east);
\end{tikzpicture}
\caption{The monolithic inverse-design procedure. The pre-solve, station
selection and degree-of-freedom check all happen once, before any Newton
iteration; the loop then solves one square system per step with no outer
optimiser. The realisability gate reports rather than blocks, because
realisability is an inviscid-consistent quantity (\S\ref{sec:presolve}) and a
viscous target can be reachable while failing it. Algorithms~\ref{alg:monolithic}
and~\ref{alg:qrpivot} state the two shaded stages precisely.}
\label{fig:pipeline}
\end{figure}

Algorithm~\ref{alg:monolithic} states the full monolithic iteration, target
generation through release-and-verify, as pseudocode, making explicit where
each of \S\ref{sec:newton-system}--\S\ref{sec:presolve}'s mechanisms (FD
Jacobian assembly, joint under-relaxation, the trust region on $\Delta A$,
the forced-transition shim) sits inside the loop. It is a direct pseudocode
rendering of
\pth{src/cins/solver/newton.py::solve_inverse} and
\pth{src/cins/benchmarks/pipeline.py::run_pipeline}, not an idealisation
of it. Algorithm~\ref{alg:qrpivot} states the QR-pivoted station-selection
procedure identified as the identifiability guard in
\S\ref{sec:ablations} (\pth{src/cins/benchmarks/pipeline.py}, station-
selection step), reusing the T4 pre-solve sensitivity matrix rather than
building a second one.

\begin{algorithm}
\caption{Monolithic CST--Newton inverse iteration (\S\ref{sec:newton-system}--\S\ref{sec:presolve})}
\label{alg:monolithic}
\begin{algorithmic}[1]
\Require target airfoil (for $A^\ast$ generation) or externally supplied $Cp_{target}$; Bernstein order $n$; \texttt{le\_treatment}; \texttt{transition.mode}
\State Fit CST to reference geometry $\rightarrow A^\ast$ (\S\ref{sec:cst-linearity}); build $\partial\zeta/\partial A_i = C(\psi)S_i(\psi)$ once, cache (Eq.~\eqref{eq:dsurface-dA})
\State Generate target: direct-solve at $A^\ast$ (forced transition if \texttt{mode=forced}, \S\ref{sec:transition}) $\rightarrow Cp_{target}^{\text{full}}$
\State Perturb or randomize $A^\ast \rightarrow A_0$ (T8 \texttt{init} factor); assemble $G,b$ from \S\ref{sec:constraints}
\If{\texttt{init = presolve}}
  \For{2 passes} \Comment{T4 pre-solve, \S\ref{sec:presolve}}
    \State Build sensitivity matrix $M$ at current $A_0$ (finite-difference or Beta-function analytic)
    \State Solve KKT system Eq.~\eqref{eq:kkt} for $A_0$; record realisability, model\_gap (ADR-0004, two \emph{distinct} metrics)
  \EndFor
\EndIf
\State Select $M$ target stations by QR pivoting (Algorithm~\ref{alg:qrpivot}) or evenly-spaced (T8 ablation)
\State \textbf{Assert} $M+K=n_{A,\text{free}}+n_\alpha$ (Eq.~\eqref{eq:dof}, FM-1 guard) \Comment{raises before any Newton step if violated}
\State Direct-solve flow at $A_0$ (converges $U_0$) \Comment{initial point for the extended Newton loop}
\State $k \gets 0$
\While{$\|R\|,\|T\|,\|G\| >$ \texttt{rtol} and $k <$ \texttt{max\_iter}} \Comment{D-1 per-block residual norms}
  \State Assemble $\partial R/\partial U$, $\partial R/\partial\alpha$ from mfoil (analytic, unchanged)
  \State Assemble $\partial R/\partial A$ by central finite differences over $A$ only (\S\ref{sec:derivatives}), $O(n_A)$ residual evaluations
  \State Assemble $\partial T/\partial U$ (Cp-selection map), $\partial G/\partial A = G$ (constant)
  \State Solve linearised system $J\,\delta \gets -[R;T-Cp_{target};G\cdot A-b]$ (Eq.~\eqref{eq:jacobian-blocks})
  \State Compute joint under-relaxation $\omega\in(0,1]$: the smaller of mfoil's own $U$-limiter (Hk, $\theta$/$\delta^*$ bounds) and the CST trust region $\|\delta_A\|_\infty \leq$ \texttt{a\_trust\_radius} (\S\ref{sec:newton-system})
  \State $U \gets U + \omega\,\delta_U$; $A \gets A + \omega\,\delta_A$; $\alpha \gets \alpha + \omega\,\delta_\alpha$ (if free)
  \State Rebuild geometry from updated $A$ (apply\_geometry, \S\ref{sec:cst-linearity}); record D-2..D-6 diagnostics
  \State $k \gets k+1$
\EndWhile
\State \textbf{Release-and-verify} (\S\ref{sec:transition}): un-force transition; direct-solve at converged $A$; assert $\Delta c_l, \Delta c_d$ against the natural-mode target within gate
\State \Return $A$, iteration count $k$, residual history, release-and-verify record
\end{algorithmic}
\end{algorithm}

\begin{algorithm}
\caption{QR-pivoted target-station selection (\S\ref{sec:ablations})}
\label{alg:qrpivot}
\begin{algorithmic}[1]
\Require sensitivity matrix $M\in\mathbb{R}^{N_{cand}\times n_A}$ from the T4 pre-solve (\S\ref{sec:presolve}); free-coefficient index set $\mathcal{F}$; required station count $M_{req} = n_{A,\text{free}} - K + n_\alpha$
\State Restrict candidates to nodes at or aft of the prescribed-LE fraction (FM-3 mitigation, \S\ref{sec:fm})
\State Form the candidate submap $\tilde M \gets M[\text{candidates},\, \mathcal{F}]$
\State Column-pivoted QR: $\tilde M^\top P = QR$, giving a pivot ordering $P$ by decreasing contribution to $\tilde M^\top$'s numerical rank
\State Select the first $M_{req}$ pivoted candidate indices $\rightarrow$ target stations
\State \Return sorted station index set; report $\mathrm{cond}(\tilde M[\text{stations},:])$ (D-3-class diagnostic, \emph{distinct} from $\mathrm{cond}(J)$, ADR-0004)
\end{algorithmic}
\end{algorithm}

\subsection{Forced-transition onset closure}
\label{sec:transition}
% <- ADR-0003 (revised mechanism); FM-4, mfoil's (and XFOIL's) $e^n$ transition closure is only $C^0$ in the design
variables: the transition-station equation itself performs an internal
sub-Newton solve for the point $x_t$ where the marched amplification factor
crosses $n_{crit}$, and that location moves discontinuously as geometry
changes cross a threshold. A coupled Newton iteration needs a continuous
residual; a kink stalls or chatters it (dossier~\S4, FM-4). The mitigation
prescribed by the dossier is to fix (trip) transition during the inverse
iterations and release it for a verification solve afterward. mfoil exposes
no forced-trip API, so this is implemented as an adapter-level shim
(ADR-0003) rather than a vendor-code edit.

The first implementation attempt froze the laminar/turbulent node flags
(\texttt{M.vsol.turb}) and no-op'd mfoil's own \texttt{update\_transition},
but left the natural \texttt{residual\_transition} equation untouched at the
frozen boundary station. This was \emph{insufficient}. Empirically (NACA
2412, $\alpha=2^\circ$, $Re=10^6$, trip 5\%/5\%), the physics moved correctly
under Newton pressure ($c_d$: 0.00578 $\to$ 0.00741), but the solve never
converged. \texttt{Rnorm} stalled at $\approx12.6$ for the full 50-iteration
budget, because the frozen station equation still solves internally for
amplification $= n_{crit}$, which has no solution once transition is pinned
somewhere amplification is nowhere near $n_{crit}$ (ADR-0003, ``Revision'').
The corrected mechanism, \texttt{\_residual\_transition\_forced}, replaces
the boundary-station residual with an XFOIL-style turbulence-\emph{onset}
closure: the momentum and shape-parameter rows use the ordinary turbulent
\texttt{residual\_station} closure directly across the whole station (no
laminar sub-leg, no station-local $x_t$), and the shear-lag row is replaced
with the same onset condition mfoil's own \texttt{update\_transition} uses to
initialize a newly-turbulent node, $U_2[2] = c_{ttr}(U_2)$, which depends
only on the downstream state (ADR-0003). With this closure, the shim
converges cleanly: trip 5\%/5\% reaches \texttt{glob.conv=True},
$\mathrm{Rnorm}<10^{-10}$, $c_d=0.01125$ (strictly $>$ natural
$c_d=0.005778$, as physically expected for larger turbulent wetted area).
Trip 30\%/30\% gives $c_d=0.00788$, strictly between the 5\%-trip and natural
values, monotonic in trip location. Releasing transition and re-solving
from the converged forced state reproduces the pinned natural baseline to
within its stated tolerance ($c_l=0.449351\pm10^{-3}$,
$c_d=0.005778\pm2\times10^{-4}$). These figures are from the shim's own
verification run (ADR-0003) and differ in the fourth decimal from the T7
archive's release-and-verify values quoted in \S\ref{sec:falsifiable}
($c_l=0.449091$, $c_d=0.005774$), which is a different run at a different
paneling, not a restatement of the same one. Rebuilding the global system after
convergence with the shim still active reproduces $\|R\|<10^{-8}$: the
converged state is a genuine fixed point of the forced residual, not an
artifact of the step sequence (ADR-0003, ``Consequences''). Every T7/T8
result reported in \S\ref{sec:results} therefore states its transition mode
explicitly, and release-and-verify in natural-transition direct mode is part
of the T7 protocol (\S\ref{sec:falsifiable}), not an optional extra.

Figure~\ref{fig:viscous} shows the state this closure operates on, and is
the reason a coupled solver was required rather than an inviscid one with a
boundary-layer post-process. The inverse system of
\S\ref{sec:newton-system} is appended to mfoil's \emph{viscous} global
system, so the six distributions in the figure are unknowns inside the same
Newton solve that recovers the geometry, not quantities computed afterwards
from a converged pressure field. Panel~(f) makes the coupling explicit: the
displacement thickness $\delta^*$ displaces the effective body the inviscid
half sees, and it is through that term that a geometry perturbation reaches
the boundary layer and back. The trip locations are marked on each panel;
the skin-friction rise at the trip in panel~(c) and the corresponding drop
in shape factor in panel~(d) are the closure described above taking effect.

\begin{figure}[htbp]
\centering
\includegraphics[width=\textwidth]{fig_viscous_solution.png}
\caption{The coupled viscous solution the inverse system is embedded in,
for NACA 2412 at $\alpha=2^\circ$, $\Rey=10^6$, with transition forced at 5
percent chord on both surfaces (dotted lines).
(a)--(b)~momentum and displacement thickness; (c)~skin friction, rising at
the trip; (d)~kinematic shape factor; (e)~edge velocity, the quantity the
inviscid half exchanges with the boundary layer; (f)~the displacement
surface $\zeta\pm\delta^*$, which is the coupling term itself. All six are
state variables of the same global system to which the CST coefficients are
appended.}
\label{fig:viscous}
\end{figure}

\subsection{Analytic pre-solve and realisability}
\label{sec:presolve}
% <- T4; ADR-0004 realisability semantics

Two jobs are done by a single linear pre-solve before any Newton iteration
runs, and the distinction between them is load-bearing for how results are
reported (ADR-0004). Thin-airfoil theory is linear in camber slope and
thickness, and CST camber/thickness are linear in $A$, so
$Cp_{target}\approx M\cdot A$ for a sensitivity matrix $M$ built either
analytically (the Joseph--Mohan Beta-function kernels) or, as implemented,
numerically from mfoil's own linear-vorticity panel influence coefficients
at fixed geometry. Solving the constrained least-squares problem with the
\S\ref{sec:constraints} rows as linear equality constraints,
\begin{equation}
  \begin{bmatrix} M^\top M & G^\top \\ G & 0 \end{bmatrix}
  \begin{bmatrix} A \\ \lambda \end{bmatrix}
  = \begin{bmatrix} M^\top Cp_{target} \\ b \end{bmatrix},
  \label{eq:kkt}
\end{equation}
gives (1) an initial guess placing Newton inside its basin
of attraction: the LE region is exactly where Newton wanders, and mfoil's
own documentation warns the coupled solver can fail from a poor initial
guess; and (2) a \textbf{realisability} metric,
$\|M A^\ast - Cp_{target}\|/\|Cp_{target}\|$, that flags targets outside the
CST-representable manifold \emph{before} a Newton solve is attempted.

An adversarial T4/T5/T7 review found these two jobs' outputs had been
conflated into one number, and ADR-0004 is the resulting decision, binding
for every number reported in this paper:

\begin{quote}
1.\ \textbf{CST-representability} (the realisability guard, above) must be
answered with \emph{model-consistent} quantities: inviscid target against
inviscid linearisation. On the T7 winning configuration this is
$3.90\times10^{-4}$; the target is
representable. 2.\ \textbf{Model gap}, how far the inviscid linearisation
sits from the viscous physics the Newton solve actually faces, is a
\emph{separate}, diagnostic-only quantity, $0.110$ on the same case. This is boundary-layer displacement
effect, which no inviscid model captures. It is \emph{not} target
unrealisability; the subsequent monolithic solve converging to
$2.75\times10^{-11}$ (\S\ref{sec:falsifiable}) is direct proof the target
was realisable despite the nonzero model gap.
\end{quote}

\textbf{Reporting rule, enforced throughout this paper:} realisability
numbers are always labelled with their model (``inviscid-consistent''), and
the conditioning of the QR-selected station submap ($\mathrm{cond}\approx90$,
\S\ref{sec:ablations}) is never quoted as the conditioning of the full
extended Newton system. These are different
matrices answering different questions; conflating them was one of the
review's confirmed findings.
